% ==============================================================================
% Fichier : front/00_page_titre.tex
% Titre   : Page de titre - Contrôle et Audit des Systèmes d'Information
% Auteur  : MINKA MI NGUIDJOÏ Thierry Emmanuel
% Version : 0.1 -- novembre 2026 (squelette)
% ==============================================================================
% Logos : Université de Yaoundé II à gauche, PSSFP à droite.
% ==============================================================================

\begin{titlepage}
\thispagestyle{empty}

% --- Fond décoratif : bandeaux bleu marine et filets or ---
\begin{tikzpicture}[remember picture, overlay]
    \fill[CouleurPrimaire]
        (current page.north west) rectangle
        ([yshift=-3.6cm] current page.north east);
    \fill[CouleurAccent]
        ([yshift=-3.6cm] current page.north west) rectangle
        ([yshift=-4.0cm] current page.north east);
    \fill[CouleurPrimaire]
        (current page.south west) rectangle
        ([yshift=2.5cm] current page.south east);
    \fill[CouleurAccent]
        ([yshift=2.5cm] current page.south west) rectangle
        ([yshift=2.9cm] current page.south east);
    % En-tête institutionnel, posé dans le bandeau supérieur
    \node[align=center, text=white, text width=12.5cm] at ([yshift=-1.8cm] current page.north) {%
        {\fontsize{11}{13}\selectfont
        \textsc{République du Cameroun}}\\[0.3ex]
        {\color{CouleurAccent!60!white}\rule{0.4\paperwidth}{0.4pt}}\\[0.6ex]
        {\fontsize{10}{12}\selectfont\textbf{\tutelle}}\\[0.2ex]
        {\fontsize{9}{11}\selectfont \institutionprimaire{} (PSSFP)}\\[0.2ex]
        {\fontsize{8}{10}\selectfont \unitepedago}};
    % Logos : Université de Yaoundé II à gauche, PSSFP à droite
    \node[fill=white, rounded corners=4pt, inner sep=2pt]
        at ([xshift=1.75cm, yshift=-1.8cm] current page.north west)
        % uy2_logo.jpeg est un bandeau (MINFI, MINESUP, UY2) : on n'en garde que le sceau de l'UY2
        {\includegraphics[height=2.6cm, trim=805 172 118 12, clip]{figures/logos/uy2_logo.jpeg}};
    \node[fill=white, rounded corners=4pt, inner sep=2pt]
        at ([xshift=-1.75cm, yshift=-1.8cm] current page.north east)
        {\includegraphics[height=2.6cm]{figures/logos/pssfp_logo.jpg}};
    % Année et version, posées dans le bandeau inférieur
    \node[align=center] at ([yshift=1.25cm] current page.south) {%
        {\color{white}\bfseries\large Année académique \anneescolaire}\\[0.3ex]
        {\color{CouleurAccent!70!white}\small
        Version 0.1 | Document pédagogique de l'UE \codeUE}};
\end{tikzpicture}

\vspace*{0.6cm}

% --- Titre principal (sur deux lignes) ---
\begin{center}
    {\color{CouleurPrimaire}\rule{0.85\textwidth}{1.2pt}}\\[1.2ex]
    {\fontsize{26}{32}\selectfont\bfseries\color{CouleurPrimaire}
        Contrôle et Audit des\\[0.4ex]
        Systèmes d'Information}\\[0.8ex]
    {\color{CouleurAccent}\rule{0.85\textwidth}{2pt}}\\[0.5ex]
    {\color{CouleurPrimaire}\rule{0.85\textwidth}{0.5pt}}
\end{center}

\vspace{0.8cm}

\begin{center}
    {\fontsize{13}{16}\selectfont\itshape\color{CouleurSecondaire}
        Évaluer le contrôle interne des finances publiques\\
        dans un environnement numérique}
\end{center}

\vspace{0.5cm}

% --- Cartouche sigle et code UE ---
\begin{center}
    \begin{tikzpicture}
        \node[draw=CouleurPrimaire, fill=CouleurPrimaire, rounded corners=6pt,
              inner xsep=14pt, inner ysep=6pt]{%
            \color{white}\bfseries\large \sigleCours{} - UE~\codeUE
        };
    \end{tikzpicture}
\end{center}

\vspace{1.0cm}

% --- Schéma de la discipline : l'évaluation du système de contrôle interne ---
\begin{center}
\begin{tikzpicture}[
    niv/.style = {
        draw=CouleurBordDef, fill=CouleurDef, rounded corners=3pt,
        minimum width=2.7cm, minimum height=0.75cm, align=center,
        font=\footnotesize\bfseries\color{CouleurPrimaire}, inner sep=2pt
    },
    arrow/.style = {-{Stealth[length=4pt]}, thick, color=CouleurSecondaire}
]
\foreach \t [count=\i] in {Mandat, Décrire, Questionner, Cheminer} {
    \node[niv] (n\i) at (3.3*\i, 0) {\t};
}
\foreach \t [count=\i from 5] in {Conception, Fonctionnement, Déficiences, Conclusion} {
    \node[niv] (n\i) at (3.3*\i - 13.2, -1.25) {\t};
}
\foreach \i [evaluate=\i as \j using int(\i+1)] in {1,2,3,5,6,7}
    \draw[arrow] (n\i) -- (n\j);
\draw[arrow] (n4.south) -- ++(0,-0.25) -| (n5.north);
\node[draw=CouleurPrimaire, fill=CouleurPrimaire, rounded corners=4pt,
      align=center, font=\footnotesize\bfseries\color{white}, inner xsep=8pt]
      (mes) at (8.25, -2.65)
      {\OUI{} / \NON{} / \NEK{} \quad$\longrightarrow$\quad
       le contrôle interne est-il efficace~?};
\draw[arrow] (n7.south) -- (n7.south |- mes.north);
\end{tikzpicture}
\end{center}

\vspace{0.6cm}

\begin{center}
    {\large\bfseries\color{CouleurPrimaire}
        \auteurcours}\\[0.5ex]
    {\normalsize\color{CouleurSecondaire}
       Enseignant Chercheur | \laboratoire{} | ENSPY}\\[0.2ex]
    {\small\color{CouleurBordInfo}
        Intervenant au PSSFP · \tutelle}\\[0.2ex]
    {\small\color{CouleurBordInfo}
        \href{mailto:maletyon@proton.me}{maletyon@proton.me}}
\end{center}

\end{titlepage}
